% 冒烟编译 ch9_c
% 《多体量子场论》中文重排本
% Xiao-Gang Wen, Quantum Field Theory of Many-Body Systems:
%   From the Origin of Sound to an Origin of Light and Electrons
%   (Oxford University Press, 2004)
% 编译：xelatex ×3
\documentclass[11pt,a4paper,twoside,openright]{ctexbook}
\usepackage{amsmath,amssymb}
\usepackage{mathrsfs}
\usepackage[margin=2.4cm]{geometry}
\usepackage{graphicx}
\usepackage{float}
\usepackage{caption}
\usepackage{needspace}
\usepackage{booktabs}
\usepackage{multirow}
\usepackage{enumitem}
\ctexset{
  chapter/name = {第,章},
  section/format = \Large\bfseries,
}
\setcounter{secnumdepth}{3}% 原书有 N.M.K.L 四级标题（如 3.3.8.1），需给 \subsubsection 编号
\setcounter{tocdepth}{1}% 目录只到节一级（与原书目录一致）
\allowdisplaybreaks
\raggedbottom
\clubpenalty=10000 \widowpenalty=10000 \predisplaypenalty=10000
\graphicspath{{figures/}}
\xeCJKDeclareCharClass{CJK}{"201C, "201D, "2018, "2019}

% 插入的空白衬页真正空白（无页眉页脚）
\makeatletter
\def\cleardoublepage{\clearpage\if@twoside\ifodd\c@page\else
\hbox{}\thispagestyle{empty}\newpage\if@twocolumn\hbox{}\newpage\fi\fi\fi}
\makeatother

% 原书小字号段落（评注/进阶话题，正文排版但字号略小）
\newenvironment{smallnote}{\par\begingroup\small}{\endgroup\par}

% 章首题词（原书各章卷首引语）
\newcommand{\epigraphbook}[2]{%
  \begin{center}\itshape #1\\[0.3em]\scshape #2\end{center}}

% 节首要点列表（原书每节开头的破折号要点）
\newenvironment{keypoints}{\begin{itemize}[leftmargin=1.6em,itemsep=0.1em]
  \setlength\parskip{0pt}}{\end{itemize}}

\begin{document}
\mainmatter
\begin{keypoints}
\item 测试要点一。
\item 测试要点二。
% ch9_c 第9章 §9.2.2尾–§9.2.7 (书页 380–392 / p395–p407)
% 块界说明：
%   【起点】§9.2.2 节标题与其要点列表前两条在书页 379（ch9_b）；书页 380（p395）
%   顶部为该要点列表的第三条，是 ch9_b 溢出到本块的内容。经协调确认：本块正文以
%   %JOIN% 起头，先续译该要点（\item，与 ch9_b 的 keypoints 环境缝合），随后在本
%   块内闭合 \end{keypoints}，再接 §9.2.2 正文。ch9_b 文件应以未闭合的 keypoints
%   收尾（末行为第二条 \item，不写 \end{keypoints}）。
%   【终点】书页 392（p407）以式 (9.2.44) 自然结束（p408 顶部的图 9.5 与
%   "The above describes..." 段属 ch9_d），无溢出，本块不设 TAIL 区。
%JOIN%
\item 非共线的 $SU(2)$ 通量把 $SU(2)$ 规范结构破缺到 $Z_2$ 规范结构。此时，所有的 $SU(2)$ 规范玻色子都获得能隙。
\end{keypoints}

正如 $U(1)$ 投影构造一样，为了在 $SU(2)$ 投影构造中得到一个稳定的平均场态，我们需要找到消去无能隙 $SU(2)$ 规范玻色子的办法。$SU(2)$ Chern--Simons 项是赋予 $SU(2)$ 规范玻色子非零能隙的一种方式。下面我们将讨论如何用 Anderson--Higgs 机制消去无能隙的 $SU(2)$ 规范玻色子。

我们知道，$U(1)$ 规范玻色子不携带自身的规范电荷，因此需要额外的带电玻色子来实现 Anderson--Higgs 机制。与之相反，$SU(2)$ 规范玻色子自身就携带非零的规范电荷，所以我们不需要额外的希格斯玻色子来实现 Anderson--\allowbreak Higgs 机制。$SU(2)$ 规范玻色子有能力杀死它们自己。

为了看清 $SU(2)$ 规范玻色子如何自杀，我们考虑一个格点 $SU(2)$ 规范理论。格点 $SU(2)$ 规范场由链变量 $U_{ij}\in SU(2)$ 给出。一阶平均场理论 (9.2.18) 就是一个 $SU(2)$ 格点规范理论。一个位形的能量是 $U_{ij}$ 的函数，即 $E(U_{ij})$。该能量在 $SU(2)$ 规范变换
\begin{equation*}
E(\tilde{U}_{ij})=E(U_{ij}),\qquad \tilde{U}_{ij}=W_i(U_{ij})W_j,\qquad W_i\in SU(2)
\tag{9.2.19}
\end{equation*}
下不变。

为了理解格点 $SU(2)$ 规范涨落的动力学，我们把 $U_{ij}$ 写成 $U_{ij}=\bar{U}_{ij}\mathrm{e}^{\mathrm{i}a_{ij}^{l}\tau^{l}}$，其中链接上的 $2\times 2$ 矩阵 $a_{ij}$ 描写规范涨落。于是能量可以写成 $E(\bar{U}_{ij},\mathrm{e}^{\mathrm{i}a_{ij}^{l}\tau^{l}})$。要判断 $SU(2)$ 规范涨落是否获得能隙，我们需要考察在 $a_{ij}^{l}$ 小的极限下，$E(\bar{U}_{ij},\mathrm{e}^{\mathrm{i}a_{ij}^{l}\tau^{l}})$ 中是否含有质量项 $(a_{ij}^{l})^{2}$。

为了理解平均场拟设 $\bar{U}_{ij}$ 如何影响规范涨落的动力学，引入平均场解的下列圈变量（loop variable）是方便的：
\begin{equation*}
P(C_i)=\bar{U}_{ij}\bar{U}_{jk}\dots\bar{U}_{li}
\tag{9.2.20}
\end{equation*}
这里 $P(C_i)$ 称为通过回路 $C_i$（由 $i\to j\to k\to\dots\to l\to i$ 给出，基点（base point）为 $i$）的 $SU(2)$ 通量，我们也称之为 $SU(2)$ 通量算符。在连续极限下，该圈变量对应于规范场强。在规范变换下，$P(C_i)$ 按如下方式变换：
\begin{equation*}
P(C_i)=W_iP(C_i)W_i
\tag{9.2.21}
\end{equation*}

我们注意到，$SU(2)$ 通量具有 $P(C)=\chi^{0}(C)\tau^{0}+\mathrm{i}\chi^{l}(C)\tau^{l}$ 的形式。因此，当 $\chi^{l}\neq 0$ 时，$SU(2)$ 通量在 $SU(2)$ 空间中有一个方向，该方向由 $\chi^{l}$ 指示。由式 (9.2.21) 可见，局域 $SU(2)$ 规范变换会转动 $SU(2)$ 通量的方向。由于具有不同基点的回路的 $SU(2)$ 通量方向可以被局域 $SU(2)$ 规范变换独立地转动，直接比较不同基点的 $SU(2)$ 通量方向是没有意义的。但是，比较具有同一基点的回路的 $SU(2)$ 通量方向却十分有意义。基于通过具有同一基点的回路的 $SU(2)$ 通量，我们可以把不同的 $SU(2)$ 通量位形分成以下三类：平凡 $SU(2)$ 通量，其中所有 $P(C)\propto\tau^{0}$；共线 $SU(2)$ 通量，其中所有 $SU(2)$ 通量都指向同一方向；非共线 $SU(2)$ 通量，其中具有同一基点的各回路的 $SU(2)$ 通量方向不同。

首先，让我们考虑对所有回路都具有平凡 $SU(2)$ 通量的拟设 $\bar{U}_{ij}$。$SU(2)$ 通量在 $SU(2)$ 规范变换下不变。我们可以（通过施行规范变换）选出一个平均场拟设，使得所有 $\bar{U}_{ij}\propto\tau^{0}$。此时，能量的规范不变性意味着
\begin{equation*}
E(\bar{U}_{ij},\mathrm{e}^{\mathrm{i}a_{ij}^{l}\tau^{l}})=E(\bar{U}_{ij},\mathrm{e}^{\mathrm{i}\theta_i^{l}\tau^{l}}\mathrm{e}^{\mathrm{i}a_{ij}^{l}\tau^{l}}\mathrm{e}^{-\mathrm{i}\theta_j^{l}\tau^{l}}).
\tag{9.2.22}
\end{equation*}
于是，在 $E$ 的展开中，质量项 $(a_{ij}^{1})^{2}$、$(a_{ij}^{2})^{2}$ 和 $(a_{ij}^{3})^{2}$ 一个都不允许出现。\footnote{在规范变换 $\mathrm{e}^{\mathrm{i}\theta_i^{1}\tau^{1}}$ 下，$a_{ij}^{1}$ 按 $a_{ij}^{\prime 1}=a_{ij}^{1}+\theta_i^{1}-\theta_j^{1}$（译注：原文左端无撇号）变换。质量项 $(a_{ij}^{1})^{2}$ 在这样的变换下不是不变的，因而不被允许。}因此，$SU(2)$ 规范涨落是无能隙的，并出现在低能处。由于拟设 $\bar{U}_{ij}\propto\tau^{0}$ 在整体 $SU(2)$ 规范变换下不变，我们说 $SU(2)$ 规范结构未破缺。

第二，让我们假设 $SU(2)$ 通量是共线的。这意思是，具有同一基点的不同回路的 $SU(2)$ 通量都指向同一方向。不过，具有不同基点的回路的 $SU(2)$ 通量仍可能指向不同方向（即使对共线 $SU(2)$ 通量也是如此）。利用局域 $SU(2)$ 规范变换，我们总能把不同基点的 $SU(2)$ 通量转到同一方向，并且可以选这个方向为 $\tau^{3}$ 方向。此时，所有 $SU(2)$ 通量都具有 $P(C)\propto\chi^{0}(C)+\mathrm{i}\chi^{3}(C)\tau^{3}$ 的形式。我们可以（通过施行 $SU(2)$ 规范变换）选出一个平均场拟设，使得所有 $\bar{U}_{ij}$ 都具有 $\mathrm{i}\,\mathrm{e}^{\mathrm{i}\phi_{ij}\tau^{3}}$ 的形式。由于该拟设在整体 $U(1)$ 规范变换 $\mathrm{e}^{\mathrm{i}\theta\tau^{3}}$ 下不变，而在 $\mathrm{e}^{\mathrm{i}\theta\tau^{1,2}}$ 下不然，我们说 $SU(2)$ 规范结构破缺到了 $U(1)$ 规范结构。能量的规范不变性意味着
\begin{equation*}
E(\bar{U}_{ij},\mathrm{e}^{\mathrm{i}a_{ij}^{l}\tau^{l}})=E(\bar{U}_{ij},\mathrm{e}^{\mathrm{i}\theta_i\tau^{3}}\mathrm{e}^{\mathrm{i}a_{ij}^{l}\tau^{l}}\mathrm{e}^{-\mathrm{i}\theta_j\tau^{3}}).
\tag{9.2.23}
\end{equation*}
当 $a_{ij}^{1,2}=0$ 时，上式约化为
\begin{equation*}
E(\bar{U}_{ij},\mathrm{e}^{\mathrm{i}a_{ij}^{3}\tau^{3}})=E(\bar{U}_{ij},\mathrm{e}^{\mathrm{i}(a_{ij}^{3}+\theta_i-\theta_j)\tau^{3}}).
\tag{9.2.24}
\end{equation*}
我们发现，质量项 $(a^{3})^{2}$ 与式 (9.2.24) 不相容，因此至少规范玻色子 $a^{3}$ 是无质量的。那么 $a^{1}$ 和 $a^{2}$ 规范玻色子又如何？令 $P_A(\pmb{i})$ 为通过基点为 $\pmb{i}$ 的回路的 $SU(2)$ 通量。如果我们假设一切能够出现在能量函数中的规范不变项都确实出现在能量函数中，则 $E(U_{ij})$ 将含有下列规范不变项：
\begin{equation*}
E=a\,\mathrm{Tr}[P_A(\pmb{i})U_{\pmb{i},\pmb{i}+\pmb{x}}P_A(\pmb{i}+\pmb{x})U_{\pmb{i}+\pmb{x},\pmb{i}}]+\dots
\tag{9.2.25}
\end{equation*}
如果把 $U_{\pmb{i},\pmb{i}+\pmb{x}}$ 写成 $\chi\,\mathrm{e}^{\mathrm{i}\phi_{ij}\tau^{3}}\mathrm{e}^{\mathrm{i}a_{x}^{l}\tau^{l}}$，利用 $U_{\pmb{i},\pmb{i}+\pmb{x}}=-U_{\pmb{i}+\pmb{x},\pmb{i}}^{\dagger}$ 这一事实（见式 (9.2.13)），并展开到 $(a_{x}^{l})^{2}$ 阶，则式 (9.2.25) 变为
\begin{equation*}
E=-\frac{1}{2}a\chi^{2}\mathrm{Tr}([P_A,a_{x}^{l}\tau^{l}]^{2})+\dots
\tag{9.2.26}
\end{equation*}
由式 (9.2.26) 我们看到，若 $P_A\propto\tau^{3}$，则 $a^{1}$ 与 $a^{2}$ 的质量项会被产生。

总结起来，我们发现：若 $SU(2)$ 通量是共线的，则拟设在 $U(1)$ 转动 $\mathrm{e}^{\mathrm{i}\theta\pmb{n}\cdot\pmb{\tau}}$（$\pmb{n}$ 为 $SU(2)$ 通量的方向）下不变。于是 $SU(2)$ 规范结构破缺到 $U(1)$ 规范结构，相应的平均场态将具有无能隙的 $U(1)$ 规范涨落。

第三，我们考虑 $SU(2)$ 通量非共线的情形。在上面我们已经证明，一个 $SU(2)$ 通量 $P_A$ 可以诱生 $\mathrm{Tr}([P_A,a_{x}^{l}\tau^{l}]^{2})$ 形式的质量项。对于非共线的 $SU(2)$ 通量位形，可以有两个指向不同方向的 $SU(2)$ 通量 $P_A$ 与 $P_B$，质量项还将含有 $\mathrm{Tr}([P_B,a_{x}^{l}\tau^{l}]^{2})$ 项。这样一来，$SU(2)$ 规范场所有的质量项，即 $(a_{ij}^{1})^{2}$、$(a_{ij}^{2})^{2}$ 和 $(a_{ij}^{3})^{2}$，都会被产生。所有的 $SU(2)$ 规范玻色子都将获得能隙（见 3.7.5 节），拟设所描写的平均场态将是一个稳定态。由于该拟设在 $Z_2$ 变换 $W_i=-\tau^{0}$ 下不变，而在其他更一般的整体 $SU(2)$ 规范变换下不然，非共线 $SU(2)$ 通量把 $SU(2)$ 规范结构破缺到 $Z_2$ 规范结构。于是我们可以猜测，低能有效理论是一个 $Z_2$ 规范理论。在 9.2.4 节和 9.3 节中，我们将研究具有非共线 $SU(2)$ 通量的态的低能性质。我们将证明，这类态的低能性质——例如 $Z_2$ 涡旋的存在和基态简并——确实与 $Z_2$ 规范理论的相同。因此，我们将把具有非共线 $SU(2)$ 通量的平均场态称为平均场 $Z_2$ 态。

在平均场 $Z_2$ 态中，所有规范涨落都有能隙。这些涨落只能在自旋子之间传递短程相互作用。低能自旋子相互作用很弱，表现得像自由自旋子。因此，把平均场涨落包括进来并不会定性地改变平均场态的性质。平均场态在低能下是稳定的。

一个稳定的平均场自旋液体态暗示着一个真实的物理自旋液体的存在。稳定平均场态的物理性质适用于该物理自旋液体。如果我们相信这两条论断，那么就可以通过研究相应的稳定平均场态来研究物理自旋液体的性质。由于在平均场 $Z_2$ 态中自旋子不被禁闭，由平均场 $Z_2$ 态得到的物理自旋液体含有中性的自旋-1/2 费米子作为其激发。意识到自旋模型是一个纯玻色模型之后，这是一个十分惊人的结果。

\subsection*{问题 9.2.3}

(a) 证明拟设 (9.2.5) 是一个 $SU(2)$ 共线态。

(b) 找出把 $U_{ij}$ 变换到 $\mathrm{e}^{\mathrm{i}\phi_{ij}\tau^{3}}$ 形式的 $SU(2)$ 规范变换。（提示：试试 $W_i=(\mathrm{i}\tau^{3})^{i_x+i_y}$。）

\subsection*{问题 9.2.4}

(a) 证明由式 (9.1.22) 得到的手征自旋拟设 $U_{ij}$ 在整体 $SU(2)$ 规范变换 $U_{ij}=WU_{ij}W^{\dagger}$ 下不变。

(b) 证明 $a_0^{l}=0$ 的手征自旋拟设 $U_{ij}$ 满足约束 (9.2.8)。

(c) 证明式 (9.1.22) 所描写的拟设与下列平移不变拟设
\begin{equation*}
\begin{aligned}
u_{\pmb{i},\pmb{i}+\pmb{x}} &= -\chi\tau^{3}-\chi\tau^{1}, &\qquad u_{\pmb{i},\pmb{i}+\pmb{y}} &= -\chi\tau^{3}+\chi\tau^{1},\\
u_{\pmb{i},\pmb{i}+\pmb{x}+\pmb{y}} &= \eta\tau^{2}, &\qquad u_{\pmb{i},\pmb{i}-\pmb{x}+\pmb{y}} &= -\eta\tau^{2},\\
a_0^{l} &= 0. &&
\end{aligned}
\end{equation*}
规范等价。用 $f$ 费米子的语言（见式 (9.2.4) 与 (9.2.7)），上述拟设描写一个 $d_{x^{2}-y^{2}}+\mathrm{i}d_{xy}$“超导”态。

由于 $SU(2)$ 规范结构未破缺，手征自旋态的低能有效理论实际上是一个 $SU(2)$ Chern--Simons 理论（level 为 1）。

\subsection{由平移不变拟设得到的自旋液体}

\begin{keypoints}
\item 均匀 RVB 态、$\pi$ 通量态与交错通量态（staggered-flux state）。
\item 投影后的物理波函数的对称性。
\item 低能 $SU(2)$ 或 $U(1)$ 规范涨落。
\end{keypoints}

作为 $SU(2)$ 投影构造的一个应用，我们将研究由平移不变拟设
\begin{equation*}
U_{\pmb{i}+\pmb{l},\pmb{j}+\pmb{l}}=U_{\pmb{i}\pmb{j}},\qquad a_0^{l}(\pmb{i})=a_0^{l},
\end{equation*}
得到的自旋液体。这样的拟设将给出具有平移对称性的物理波函数。\footnote{这里我们要区分拟设的不变性与拟设的对称性。当拟设自身在平移下不改变时，我们说拟设具有平移不变性；当由拟设得到的物理自旋波函数具有平移对称性时，我们说拟设具有平移对称性。由于规范结构的存在，一个不具有平移不变性的拟设也可以具有平移对称性。}首先，让我们引入
\begin{equation*}
u_{ij}=\frac{3}{8}J_{ij}U_{ij}=u_{ij}^{\mu}\tau^{\mu}
\end{equation*}
其中 $u_{\pmb{l}}^{1,2,3}$ 为实数，$u_{\pmb{l}}^{0}$ 为虚数。在零阶平均场理论中，自旋子谱由哈密顿量（见式 (9.2.9)）
\begin{equation*}
H=-\sum_{\langle ij\rangle}\left(\psi_i^{\dagger}u_{ij}\psi_j+h.c.\right)+\sum_i\psi_i^{\dagger}a_0^{l}\tau^{l}\psi_i
\tag{9.2.27}
\end{equation*}
决定。在 $\pmb{k}$ 空间中，我们有
\begin{equation*}
H=-\sum_{\pmb{k}}\psi_{\pmb{k}}^{\dagger}(u^{\mu}(\pmb{k})-a_0^{\mu})\tau^{\mu}\psi_{\pmb{k}}
\end{equation*}
其中 $\mu=0,1,2,3$（译注：原文此处印作“0, 1, 2,”），
\begin{equation*}
u^{\mu}(\pmb{k})=\sum_{\pmb{l}}u_{\pmb{i},\pmb{i}+\pmb{l}}^{\mu}\mathrm{e}^{\mathrm{i}\pmb{l}\cdot\pmb{k}},
\end{equation*}
$a_0^{0}=0$，而 $N$ 为格点总数。费米子谱有两支，由
\begin{align*}
E_{\pm}(\pmb{k})&=u^{0}(\pmb{k})\pm E_0(\pmb{k})\\
E_0(\pmb{k})&=\sqrt{\sum_{\pmb{l}}(u^{l}(\pmb{k})-a_0^{l})^{2}}
\tag{9.2.28}
\end{align*}
给出。约束可以由 $\partial E_{\mathrm{ground}}/\partial a_0^{l}=0$ 得到，其形式为
\begin{equation*}
N\langle\psi_i^{\dagger}\tau^{l}\psi_i\rangle=\sum_{\pmb{k},E_-(\pmb{k})<0}\frac{u^{l}(\pmb{k})-a_0^{l}}{E_0(\pmb{k})}-\sum_{\pmb{k},E_+(\pmb{k})<0}\frac{u^{l}(\pmb{k})-a_0^{l}}{E_0(\pmb{k})}=0
\tag{9.2.29}
\end{equation*}
由此可以定出 $a_0^{l}$，$l=1,2,3$。

让我们尝试下列简单拟设\footnote{该拟设实际上就是前面引入的 $d$ 波拟设 (9.2.5)。}（Baskaran et al., 1987; Affleck and Marston, 1988; Kotliar and Liu, 1988）：
\begin{equation*}
a_0^{l}=0,\qquad u_{\pmb{i},\pmb{i}+\pmb{x}}=\chi\tau^{3}+\eta\tau^{1},\qquad u_{\pmb{i},\pmb{i}+\pmb{y}}=\chi\tau^{3}-\eta\tau^{1}.
\tag{9.2.30}
\end{equation*}
首先，我们要证明相应的自旋液体具有全部对称性。

为了研究相应物理自旋态的对称性，我们需要求出它的波函数 $|\Psi_{\mathrm{phy}}\rangle$。利用该拟设，可以得到平均场哈密顿量 (9.2.27) 的平均场基态 $|\Psi_{\mathrm{mean}}\rangle$。把 $|\Psi_{\mathrm{mean}}\rangle$ 投影到每格点 $\psi$ 费米子数为偶数的子空间，就得到物理波函数 $|\Psi_{\mathrm{phy}}\rangle$，即 $|\Psi_{\mathrm{phy}}\rangle=\mathcal{P}|\Psi_{\mathrm{mean}}\rangle$。

该拟设已经在这两个平移 $T_x:\,\pmb{i}\to\pmb{i}+\pmb{x}$ 与 $T_y:\,\pmb{i}\to\pmb{i}+\pmb{y}$ 下不变。拟设也在两个宇称变换 $P_x:\,x\to-x$ 与 $P_y:\,y\to-y$ 下不变。宇称变换 $P_{xy}:\,(x,y)\to(y,x)$ 把 $u_{\pmb{i},\pmb{i}+\pmb{x}}$ 变为 $u_{\pmb{i},\pmb{i}+\pmb{y}}$、把 $u_{\pmb{i},\pmb{i}+\pmb{y}}$ 变为 $u_{\pmb{i},\pmb{i}+\pmb{x}}$。变换后的拟设
\begin{equation*}
a_0^{l}=0,\qquad u_{\pmb{i},\pmb{i}+\pmb{x}}=\chi\tau^{3}-\eta\tau^{1},\qquad u_{\pmb{i},\pmb{i}+\pmb{y}}=\chi\tau^{3}+\eta\tau^{1}.
\end{equation*}
与原来的拟设不同。但是请注意，一个整体规范变换 $G_{P_{xy}}(\pmb{i})=\mathrm{i}\tau^{3}$ 把 $u_{ij}$ 变为 $G_{P_{xy}}u_{ij}G_{P_{xy}}^{\dagger}$，它诱导 $(\tau^{1},\tau^{2},\tau^{3})\to(-\tau^{1},-\tau^{2},\tau^{3})$ 的改变。因此，变换后的拟设通过规范变换 $G_{P_{xy}}$ 与原来的拟设规范等价。换句话说，$P_{xy}$ 宇称变换再接着一个规范变换 $G_{P_{xy}}$，保持拟设不变。于是，物理波函数在 $P_{xy}$ 宇称变换下不变。$90^{\circ}$ 转动 $R_{90}$ 由 $P_xP_{xy}$ 生成，故物理波函数也有 $90^{\circ}$ 转动对称性。上述拟设还有时间反演对称性，因为时间反演变换 $u_{ij}\to-u_{ij}$ 再接着规范变换 $G_T(i)=\mathrm{i}\tau^{2}$，保持拟设不变。

总而言之，该拟设在下列对称变换与规范变换的组合 $G_xT_x$、$G_yT_y$、$G_{P_x}P_x$、$G_{P_y}P_y$、$G_{P_{xy}}P_{xy}$ 与 $G_TT$ 下不变，相应的规范变换由
\begin{equation*}
\begin{array}{ccc}
G_x=\tau^{0}, & \quad G_y=\tau^{0}, & \quad G_T=\mathrm{i}\tau^{2},\\
G_{P_x}=\tau^{0}, & \quad G_{P_y}=\tau^{0}, & \quad G_{P_{xy}}=\mathrm{i}\tau^{3}.
\end{array}
\tag{9.2.31}
\end{equation*}
给出。所以，由拟设 (9.2.30) 描写的自旋液体具有全部对称性。我们把这样的态称为对称自旋液体态。

我们注意到，式 (9.2.30) 中的链变量 $u_{ij}$ 在 $\tau^{1,2,3}$ 空间中指向不同的方向，因而不是共线的。这样一来，该拟设在整体 $SU(2)$ 规范变换下并非不变。人们可能因此期待 $SU(2)$ 规范结构破缺到 $Z_2$ 规范结构，从而拟设 (9.2.30) 描写一个 $Z_2$ 自旋液体。实际上这个推理是错误的。正如我们在 9.2.2 节中指出的，要判断 $SU(2)$ 规范结构是否破缺到 $Z_2$ 规范结构，需要检查的是 $SU(2)$ \emph{通量}的共线性，而不是 $SU(2)$ 链变量的共线性。

当 $\chi=\eta$ 或 $\eta=0$ 时，所有回路的 $SU(2)$ 通量 $P_C$ 都是平凡的，即 $P_C\propto\tau^{0}$。此时 $SU(2)$ 规范结构未破缺，$SU(2)$ 规范涨落是无能隙的（在弱耦合极限下）。$\eta=0$ 所描写的自旋液体中的自旋子具有谱 $E_{\pmb{k}}=\pm 2|\chi(\cos(k_x)+\cos(k_y))|$，它有一个大费米面。我们将把这个态称为 $SU(2)$ 无能隙态（文献中称其为均匀 RVB 态）。$\chi=\eta$ 的态只在孤立的 $\pmb{k}$ 点有无能隙自旋子，自旋子谱为 $E_{\pmb{k}}=\pm 2\sqrt{\chi^{2}+\eta^{2}}\sqrt{\cos^{2}(k_x)+\cos^{2}(k_y)}$（见图 9.2）。我们将把这样的态称为 $SU(2)$ 线性态，以强调费米点附近 $E\propto|\pmb{k}|$ 的线性色散。（这种态在文献中以及在 9.1.1 节中称为 $\pi$ 通量态。）$SU(2)$ 线性态的低能有效理论由无质量狄拉克费米子（即自旋子）耦合到 $SU(2)$ 规范场所描写（见问题 9.1.4）。

经过适当的规范变换，$SU(2)$ 无能隙拟设可以改写为
\begin{equation*}
u_{\pmb{i},\pmb{i}+\pmb{x}}=\mathrm{i}\chi,\qquad\qquad u_{\pmb{i},\pmb{i}+\pmb{y}}=\mathrm{i}\chi,
\tag{9.2.32}
\end{equation*}
而 $SU(2)$ 线性拟设可以改写为
\begin{equation*}
u_{\pmb{i},\pmb{i}+\pmb{x}}=\mathrm{i}\chi,\qquad\qquad u_{\pmb{i},\pmb{i}+\pmb{y}}=\mathrm{i}(-)^{i_x}\chi.
\tag{9.2.33}
\end{equation*}
在这些形式中，由于 $u_{ij}\propto\mathrm{i}\tau^{0}$，$SU(2)$ 规范结构明显未破缺。我们可以容易地看到，所有的 $SU(2)$ 通量 $P_{ij\dots k}$ 都是平凡的。\footnote{在后面将要讨论的投影对称群（projective symmetry group）分类中，$SU(2)$ 无能隙拟设 (9.2.32) 标记为 SU2An0，$SU(2)$ 线性拟设 (9.2.33) 标记为 SU2Bn0（见式 (9.4.33)）。}这些拟设也在整体 $SU(2)$ 规范变换下不变。$SU(2)$ 无能隙拟设与 $SU(2)$ 线性拟设中的自旋子在低能下都有有限的相互作用，所以这两个平均场态都是边缘平均场态。这两个平均场态是否对应于真实的物理自旋液体，目前还不清楚。

当 $\chi\neq\eta$ 且 $\chi,\eta\neq 0$ 时，通量 $P_C$ 是非平凡的，但 $SU(2)$ 通量是共线的。此时，$SU(2)$ 规范结构破缺到 $U(1)$ 规范结构。无能隙自旋子仍然只出现在孤立的 $\pmb{k}$ 点，从自旋子谱 $E_{\pmb{k}}=\pm 2\sqrt{\chi^{2}(\cos k_x+\cos k_y)^{2}+\eta^{2}(\cos k_x-\cos k_y)^{2}}$ 可以看出这一点。我们将把这样的态称为 $U(1)$ 线性态。（这个态在文献中称为交错通量态和/或 $d$ 波配对态。）经过适当的规范变换，$U(1)$ 线性态也可以用拟设（见问题 9.2.3）
\begin{equation*}
u_{\pmb{i},\pmb{i}+\pmb{x}}=\mathrm{i}\chi-(-)^{i}\eta\tau^{3},\qquad u_{\pmb{i},\pmb{i}+\pmb{y}}=\mathrm{i}\chi+(-)^{i}\eta\tau^{3},
\tag{9.2.34}
\end{equation*}
来描写，其中 $U(1)$ 规范结构是显式的。\footnote{在投影对称群分类中，这样的态标记为 U1Cn01n（见式 (9.4.25)--(9.4.29)）。}其低能有效理论由无质量狄拉克费米子（即自旋子）耦合到 $U(1)$ 规范场所描写。同样，由于无能隙 $U(1)$ 规范场在低能诱导的有限相互作用，$U(1)$ 线性态也是一个边缘平均场态。

\subsection{由平移不变拟设得到的稳定 $Z_2$ 自旋液体}

\begin{keypoints}
\item 稳定的自旋液体态——$Z_2$ 自旋液体态。
\item $Z_2$ 自旋液体含有分数化的激发——自旋子。自旋子携带自旋-1/2、不带电荷，具有费米统计。
\end{keypoints}

在学会了构造稳定平均场态的配方之后，本节我们研究一个不破坏任何对称性的稳定平均场态。该拟设具有下列形式（Wen, 1991a）：
\begin{align*}
u_{\pmb{i},\pmb{i}+\pmb{x}}=u_{\pmb{i},\pmb{i}+\pmb{y}}=-\chi\tau^{3},\qquad & a_0^{2,3}=0,\quad a_0^{1}\neq 0,\\
u_{\pmb{i},\pmb{i}+\pmb{x}+\pmb{y}}=\eta\tau^{1}+\lambda\tau^{2},\qquad & u_{\pmb{i},\pmb{i}-\pmb{x}+\pmb{y}}=\eta\tau^{1}-\lambda\tau^{2}
\tag{9.2.35}
\end{align*}
为了证明式 (9.2.35) 描写一个稳定的平均场态，我们只需证明它产生非共线的 $SU(2)$ 通量。围绕两个具有同一基点的三角形，我们得到下列 $SU(2)$ 通量：
\begin{align*}
u_{\pmb{i},\pmb{i}+\pmb{x}}u_{\pmb{i}+\pmb{x},\pmb{i}+\pmb{y}}u_{\pmb{i}+\pmb{y},\pmb{i}}&=-(\eta\tau^{1}-\lambda\tau^{2})\chi^{2}\\
u_{\pmb{i},\pmb{i}+\pmb{y}}u_{\pmb{i}+\pmb{y},\pmb{i}-\pmb{x}}u_{\pmb{i}-\pmb{x},\pmb{i}}&=-(\eta\tau^{1}+\lambda\tau^{2})\chi^{2}
\end{align*}
我们看到，若 $\chi\eta\lambda\neq 0$，则 $SU(2)$ 通量非共线，式 (9.2.35) 中的拟设描写一个稳定的平均场态，更确切地说，一个 $Z_2$ 态。

由于拟设在平移下不变，物理自旋态在平移下也不变，于是物理波函数 $|\Psi_{\mathrm{phy}}\rangle$ 具有平移对称性。在 $P_x$ 宇称变换下，拟设变为
\begin{align*}
u_{\pmb{i},\pmb{i}+\pmb{x}}=u_{\pmb{i},\pmb{i}+\pmb{y}}=-\chi\tau^{3},\qquad & a_0^{2,3}=0,\quad a_0^{1}\neq 0,\\
u_{\pmb{i},\pmb{i}+\pmb{x}+\pmb{y}}=\eta\tau^{1}-\lambda\tau^{2},\qquad & u_{\pmb{i},\pmb{i}-\pmb{x}+\pmb{y}}=\eta\tau^{1}+\lambda\tau^{2}
\tag{9.2.36}
\end{align*}
（注意，对平移不变拟设，$P_x$ 把 $u_{\pmb{i},\pmb{i}+\pmb{x}}$ 变为 $u_{\pmb{i},\pmb{i}-\pmb{x}}=u_{\pmb{i},\pmb{i}+\pmb{x}}^{\dagger}$）。变换后的拟设在 $SU(2)$ 规范变换 $G_x(i)=\mathrm{i}\tau^{1}(-)^{i}$ 下与原来的拟设规范等价。\footnote{我们定义 $(-)^{i}\equiv(-)^{i_x+i_y}$。规范变换 $W_i=\mathrm{i}\tau^{1}$ 把式 (9.2.36) 变换为
\begin{equation*}
\begin{aligned}
u_{\pmb{i},\pmb{i}+\pmb{x}}&=u_{\pmb{i},\pmb{i}+\pmb{y}}=\chi\tau^{3},\qquad & a_0^{2,3}&=0,\quad a_0^{1}\neq 0,\\
u_{\pmb{i},\pmb{i}+\pmb{x}+\pmb{y}}&=\eta\tau^{1}+\lambda\tau^{2},\qquad & u_{\pmb{i},\pmb{i}-\pmb{x}+\pmb{y}}&=\eta\tau^{1}-\lambda\tau^{2}
\end{aligned}
\end{equation*}
然后，规范变换 $W_i=(-)^{i}$ 再把上式变为式 (9.2.35)。}尽管 $P_x$ 变换后的拟设与原来的拟设不同，两者在投影后给出同一个物理波函数。物理自旋态具有 $P_x$ 宇称对称性。

我们还注意到，若 $(a_0^{2},a_0^{3})$ 不为零，上述组合变换 $G_xP_x$ 会把它们变成 $(-a_0^{2},-a_0^{3})$。由于平均场哈密顿量具有 $P_x$ 对称性，非零 $(a_0^{2},a_0^{3})$ 的平均场基态能量满足性质 $E_{\mathrm{ground}}(a_0^{2},a_0^{3})=E_{\mathrm{ground}}(-a_0^{2},-a_0^{3})$。因此，对我们的拟设而言，$a_0^{2}=a_0^{3}=0$ 总能满足约束 $\langle\psi_i^{\dagger}\tau^{2}\psi_i\rangle\allowbreak=N_{\mathrm{site}}^{-1}\partial E_{\mathrm{ground}}/\partial a_0^{2}\allowbreak=0$ 与 $\langle\psi_i^{\dagger}\tau^{3}\psi_i\rangle\allowbreak=N_{\mathrm{site}}^{-1}\partial E_{\mathrm{ground}}/\partial a_0^{3}\allowbreak=0$。我们只需调节 $a_0^{1}$ 使 $\langle\psi_i^{\dagger}\tau^{1}\psi_i\rangle=0$。

$P_y$ 宇称变换的作用与 $P_x$ 相同，所以物理波函数也有 $P_y$ 宇称对称性。在时间反演变换下，$T$ 把 $(u_{ij};a_0)$ 变为 $(-u_{ij};-a_0)$。如果我们选 $G_T=\mathrm{i}\tau^{3}(-)^{i}$，时间反演变换再接着规范变换 $G_T$ 将保持拟设 (9.2.35) 不变。因此，$|\Psi_{\mathrm{phy}}\rangle$ 具有时间反演对称性。该拟设在 $(x,y)\to(y,x)$ 宇称 $P_{xy}$ 下不变，所以物理波函数具有 $P_{xy}$ 宇称对称性。我们看到，拟设 (9.2.35) 在下列组合变换 $G_xT_x$、$G_yT_y$、$G_{P_x}P_x$、$G_{P_y}P_y$、$G_{P_{xy}}P_{xy}$ 与 $G_TT$ 下不变，其中规范变换为
\begin{equation*}
\begin{array}{ccc}
G_x=\tau^{0}, & \quad G_y=\tau^{0}, & \quad G_T=\mathrm{i}(-)^{i}\tau^{3},\\
G_{P_x}=\tau^{0}, & \quad G_{P_y}=\tau^{0}, & \quad G_{P_{xy}}=\mathrm{i}\tau^{3}
\end{array}
\tag{9.2.37}
\end{equation*}
该拟设描写一个具有全部对称性的自旋液体。我们把这样的态称为 $Z_2$ 对称自旋液体态。

$Z_2$ 态中的自旋子谱由
\begin{align*}
E_{\pm}(\pmb{k})&=\pm\sqrt{\epsilon_1^{2}+\epsilon_2^{2}+\epsilon_3^{2}}\\
\epsilon_1&=2\chi(\cos k_x+\cos k_y)\\
\epsilon_2&=2\eta[\cos(k_x+k_y)+\cos(k_x-k_y)]+a_0^{1}\\
\epsilon_3&=2\lambda[\cos(k_x+k_y)-\cos(k_x-k_y)]
\end{align*}
给出。我们发现自旋子是完全有能隙的。我们将把态 (9.2.35) 称为 $Z_2$ 有能隙自旋液体。\footnote{这里我们想指出，在 9.4.3 节将要讨论的投影对称群分类中，态 (9.2.35) 所标记的投影对称群为 Z2Axx0z。}$Z_2$ 有能隙自旋液体对应于 Kivelson et al. (1987) 以及 Rokhsar and Kivelson (1988) 提出的短程共振价键（sRVB）态。

由于式 (9.2.35) 是一个稳定拟设，规范涨落全都有能隙，只在自旋子之间产生短程相互作用。由于自旋子有能隙，$Z_2$ 有能隙态的低能激发对应于自旋子的稀薄气体。由于相互作用是短程的，这些自旋子表现得像自由费米子。因此，$Z_2$ 有能隙态中的激发是自由的中性自旋-1/2 费米子。

%FIG{9.4}: {沿粗线所示的链接翻转 $U_{ij}$ 的符号，就产生一个 $Z_2$ 涡旋。这里 $\Theta_{ij}=1$，唯有粗线所示的链接上 $\Theta_{ij}=-1$。}

\subsection{$Z_2$ 自旋液体中的 $Z_2$ 涡旋}

\begin{keypoints}
\item $Z_2$ 自旋液体态含有一种 $Z_2$ 涡旋激发，它既不携带自旋也不携带电荷。
\item 把一个 $Z_2$ 涡旋绑定到自旋子上，会使自旋子的统计性从玻色性变为费米性，或从费米性变为玻色性。
\end{keypoints}

除自旋子外，$Z_2$ 自旋液体态中还有另一类准粒子激发（Wen, 1991a）。这种激发在平均场理论中以拓扑孤子（topological soliton）的形式出现，对应于 $Z_2$ 规范场的 $\pi$ 通量，因此我们将把这种新激发称为 $Z_2$ 涡旋。它通过沿一串链接翻转 $u_{ij}$ 的符号产生（见图 9.4），由下列拟设描写：
\begin{equation*}
\tilde{u}_{ij}=\bar{u}_{ij}\Theta_{ij}
\tag{9.2.38}
\end{equation*}
其中 $\Theta_{ij}$ 已在图 9.4 中说明。$Z_2$ 涡旋位于弦的端点。注意，在远离 $Z_2$ 涡旋处，$\tilde{u}_{ij}$ 局域地规范等价于 $\bar{u}_{ij}$，因此符号翻转的 $u_{ij}$ 弦不耗费任何能量，也观测不到。拟设 $\tilde{u}_{ij}$ 实际上描写的是位于弦端点的一个局域激发。

由于 $u_{ij}$ 决定自旋子的跳跃幅，当自旋子绕 $Z_2$ 涡旋跳行一周时会多出一个负号。因此，对自旋子而言 $Z_2$ 涡旋表现得像一个 $\pi$ 通量。根据 7.1.2 节的讨论，我们发现把一个 $Z_2$ 涡旋绑定到一个自旋子上，会把自旋子的统计性从费米性变为玻色性。因此，$Z_2$ 自旋液体既含有玻色统计的中性自旋-1/2 激发，也含有费米统计的中性自旋-1/2 激发。

\subsection*{问题 9.2.5}

证明：对于式 (9.2.38) 的 $Z_2$ 涡旋拟设，通过回路 $C$ 的 $SU(2)$ 通量与基态拟设 $u_{ij}$ 通过同一回路的 $SU(2)$ 通量相同，除非回路 $C$ 包围着涡旋。由于平均场能量是 $SU(2)$ 通量的函数，上述结果表明符号翻转的 $u_{ij}$ 弦不耗费能量。

\subsection{稳定的无能隙 $Z_2$ 自旋液体}

\begin{keypoints}
\item 存在许多不同的稳定自旋液体态——$Z_2$ 自旋液体态。其中一些可以具有无能隙的自旋子激发。
\item 这些态具有完全相同的对称性。没有办法利用对称性来刻画这些不同的自旋液体态。
\end{keypoints}

除了式 (9.2.35) 之外，对于 $Z_2$ 对称自旋液体我们还可以写下另一个拟设（Wen, 2002c）：
\begin{equation*}
u_{\pmb{i},\pmb{i}+\pmb{x}}=\mathrm{i}\eta\tau^{0}-\chi(\tau^{3}-\tau^{1}),\qquad u_{\pmb{i},\pmb{i}+\pmb{y}}=\mathrm{i}\eta\tau^{0}-\chi(\tau^{3}+\tau^{1}),\qquad a_0^{l}=0,
\tag{9.2.39}
\end{equation*}
其中 $\chi$ 与 $\eta$ 均不为零。当 $\chi=0$ 时，上述拟设成为 $SU(2)$ 无能隙自旋液体；当 $\eta=0$ 时，成为 $SU(2)$ 线性自旋液体。为了证明式 (9.2.39) 的拟设描写一个对称自旋液体，我们必须证明它在下列对称变换与规范变换的组合 $G_xT_x$、$G_yT_y$、$G_{P_x}P_x$、$G_{P_y}P_y$、$G_{P_{xy}}P_{xy}$ 与 $G_TT$ 下不变。我们发现，如果选取规范变换为
\begin{equation*}
\begin{array}{lll}
G_x=\tau^{0}, & \quad G_y=\tau^{0}, & \quad G_T=(-)^{i}\tau^{0},\\[2pt]
G_{P_x}=\mathrm{i}(-)^{i_x}\dfrac{\tau^{1}+\tau^{3}}{\sqrt{2}}, & \quad G_{P_y}=\mathrm{i}(-)^{i_y}\dfrac{\tau^{1}-\tau^{3}}{\sqrt{2}}, & \quad G_{P_{xy}}=\mathrm{i}\tau^{3},
\end{array}
\tag{9.2.40}
\end{equation*}
则对称变换再接着相应的规范变换，将保持拟设不变。

利用时间反演对称性，我们可以证明拟设 (9.2.39) 中取零值的 $a_0^{l}$ 确实满足约束 (9.2.29)。这是因为在组合时间反演变换 $G_TT$ 下 $a_0^{l}\to-a_0^{l}$。由于 $a_0^{l}=0$ 的拟设在 $G_TT$ 下不变，非零 $a_0^{l}$ 的平均场基态能量满足 $E_{\mathrm{ground}}(a_0^{l})=E_{\mathrm{ground}}(-a_0^{l})$，于是当 $a_0^{l}=0$ 时 $\partial E_{\mathrm{ground}}/\partial a_0^{l}=0$。事实上，任何只在两个不交叠子格之间有链接的拟设（即非阻挫拟设），只要 $a_0^{l}=0$，就是时间反演对称的。对这类拟设（包括拟设 (9.2.39)），取零值的 $a_0^{l}$ 满足约束 (9.2.29)。

自旋子谱由下式给出（见图 9.10(a)）
\begin{equation*}
E_{\pm}=2\eta(\sin(k_x)+\sin(k_y))\pm 2|\chi|\sqrt{2\cos^{2}(k_x)+2\cos^{2}(k_y)}
\end{equation*}
自旋子有两个费米点和两个小的费米口袋（Fermi pocket）（当 $\eta$ 小时）。$SU(2)$ 通量是非平凡的。而且，$SU(2)$ 通量 $P_{C_1}$ 与 $P_{C_2}$ 互不对易，其中 $C_1=\pmb{i}\to\pmb{i}+\pmb{x}\to\pmb{i}+\pmb{x}+\pmb{y}\to\pmb{i}+\pmb{y}\to\pmb{i}$ 与 $C_2=\pmb{i}\to\pmb{i}+\pmb{y}\to\pmb{i}-\pmb{x}+\pmb{y}\to\pmb{i}-\pmb{x}\to\pmb{i}$ 是两个具有同一基点的回路。非共线 $SU(2)$ 通量把 $SU(2)$ 规范结构破缺到 $Z_2$ 规范结构。我们将把式 (9.2.39) 描写的自旋液体称为 $Z_2$ 无能隙自旋液体。\footnote{$Z_2$ 无能隙自旋液体是 9.4.3 节中分类的 $Z_2$ 自旋液体之一，其投影对称群标记为 Z2Ax12(12)n（见 9.4.3 节及式 (9.4.16)）。}其低能有效理论由无质量狄拉克费米子和具有小费米面的费米子耦合到 $Z_2$ 规范场所描写。由于没有无能隙规范玻色子，规范涨落只能在自旋子之间传递短程相互作用。短程相互作用在低能下是不相关的，自旋子在低能下是自由费米子。因此，平均场 $Z_2$ 无能隙态是一个稳定态。这个稳定的平均场态暗示存在第二种真实的物理自旋液体（除 9.2.4 节讨论的 $Z_2$ 有能隙自旋液体之外）。这种自旋液体的物理性质可以由平均场理论得到。我们发现，$Z_2$ 无能隙自旋液体中的激发由中性的自旋-1/2 费米型准粒子描写！

第三种 $Z_2$ 对称自旋液体可以由下列阻挫拟设得到（Balents et al., 1998; Senthil and Fisher, 2000）：
\begin{align*}
a_0^{3}\neq 0,\qquad & a_0^{1,2}=0,\\
u_{\pmb{i},\pmb{i}+\pmb{x}}=\chi\tau^{3}+\eta\tau^{1},\qquad & u_{\pmb{i},\pmb{i}+\pmb{y}}=\chi\tau^{3}-\eta\tau^{1},\\
u_{\pmb{i},\pmb{i}+\pmb{x}+\pmb{y}}=+\gamma\tau^{3},\qquad & u_{\pmb{i},\pmb{i}-\pmb{x}+\pmb{y}}=+\gamma\tau^{3}.
\tag{9.2.41}
\end{align*}
该拟设具有平移、转动、宇称与时间反演对称性。它在组合变换 $G_{x,y}T_{x,y}$、$G_{P_x,P_y,P_{xy}}P_{x,y,xy}$ 与 $G_TT$ 下不变，其中
\begin{equation*}
\begin{array}{ccc}
G_x=\tau^{0}, & \quad G_y=\tau^{0}, & \quad G_T=\mathrm{i}\tau^{2},\\
G_{P_x}=\tau^{0}, & \quad G_{P_y}=\tau^{0}, & \quad G_{P_{xy}}=\mathrm{i}\tau^{3}.
\end{array}
\tag{9.2.42}
\end{equation*}
当 $a_0^{3}\neq 0$、$\chi\neq\pm\eta$ 且 $\chi\eta\neq 0$ 时，我们发现 $a_0^{l}\tau^{l}$ 与 $SU(2)$ 通量算符——例如穿过一个方格的 $SU(2)$ 通量 $P(C_{\mathrm{square}})$——不对易。\footnote{这里 $\tau^{l}a_0^{l}(\pmb{i})$ 可以看作通过一个零尺寸回路的 $SU(2)$ 通量。}因此，该拟设把 $SU(2)$ 规范结构破缺到 $Z_2$ 规范结构。自旋子谱由下式给出（见图 9.8(a)）
\begin{align*}
E_{\pm}&=\pm\sqrt{\epsilon^{2}(\pmb{k})+\Delta^{2}(\pmb{k})}\\
\epsilon(\pmb{k})&=2\chi(\cos(k_x)+\cos(k_y))+2\gamma(\cos(k_x+k_y)+\cos(k_x-k_y))+a_0^{3}\\
\Delta(\pmb{k})&=2\eta(\cos(k_x)-\cos(k_y))+a_0^{3}
\end{align*}
它只在该四个 $\pmb{k}$ 点无能隙，且具有线性色散。因此，式 (9.2.41) 描写的自旋液体是一个 $Z_2$ 线性自旋液体。\footnote{$Z_2$ 线性自旋液体由投影对称群 Z2A0032 或等价的 Z2A0013 描写（见 9.4.3 节）。}

$Z_2$ 有能隙拟设 (9.2.35)、$Z_2$ 无能隙拟设 (9.2.39) 与 $Z_2$ 线性拟设 (9.2.41) 在投影后给出三个自旋液体态。这三个态具有完全相同的对称性，所以我们不能用对称性来区分它们。为了找到一组新的量子数来刻画这三种不同的自旋液体，我们注意到：虽然三种 $Z_2$ 自旋液体具有相同的对称性，但它们的拟设分别在同样的对称平移之后跟着\emph{不同的}规范变换下不变（见式 (9.2.37)、(9.2.40) 与 (9.2.42)）。因此，三种自旋液体的平均场拟设的不变群各不相同，我们可以用这样的不变群来刻画不同的自旋液体。更详细和系统的讨论将在 9.4 节给出。

\subsection*{问题 9.2.6}

证明：如果规范变换由式 (9.2.40) 给出，则式 (9.2.39) 中的拟设在组合变换 $G_xT_x$、$G_yT_y$、$G_{P_x}P_x$、$G_{P_y}P_y$、$G_{P_{xy}}P_{xy}$ 与 $G_TT$ 下不变。

\subsection*{问题 9.2.7}

证明手征自旋态具有平移对称性和 $90^{\circ}$ 转动对称性。证明手征自旋态还具有组合对称性 $TP_x$、$TP_y$ 与 $TP_{xy}$。

\subsection{评注：平均场理论中的时间反演变换}

在时间反演变换（或自旋反转变换）下，$\pmb{S}_i\to-\pmb{S}_i$。由于本章只考虑自旋旋转不变的态，为方便起见，我们把 $T$ 变换定义为时间反演变换与自旋旋转变换的组合：$(S^{x},S^{y},S^{z})\to(S^{x},-S^{y},S^{z})$。我们将不严格地把它称为时间反演变换。

这样的变换可以由
\begin{equation*}
\psi\to\psi'=-\mathrm{i}\tau^{2}\psi^{*}
\tag{9.2.43}
\end{equation*}
生成。让我们用一个算符 $\mathcal{T}$ 来表示上述变换，即 $\mathcal{T}\psi\mathcal{T}^{-1}=-\mathrm{i}\tau^{2}\psi^{*}$。注意 $\mathcal{T}$ 把 $\mathrm{i}$ 变成 $-\mathrm{i}$，即 $\mathcal{T}^{-1}\mathrm{i}\mathcal{T}=-\mathrm{i}$。用 $f$ 费米子 (9.2.6) 表示，式 (9.2.43) 变为
\begin{equation*}
f\to\mathcal{T}f\mathcal{T}^{-1}=-\mathrm{i}\tau^{2}f^{*}
\end{equation*}
由 $\pmb{S}=\frac{1}{2}f^{\dagger}\pmb{\sigma}f$，我们发现 $\mathcal{T}$ 把 $\pmb{S}$ 变为
\begin{equation*}
\mathcal{T}\pmb{S}\mathcal{T}^{-1}=-\frac{1}{2}f^{\top}\pmb{\sigma}f^{*}=\frac{1}{2}f^{\dagger}\pmb{\sigma}^{\top}f=(S^{x},-S^{y},S^{z})
\end{equation*}
在平均场理论中，式 (9.2.43) 诱导平均场拟设的下列变换：
\begin{equation*}
\begin{aligned}
U_{ij}&\to U'_{ij}=(-\mathrm{i}\tau^{2})U_{ij}^{*}(\mathrm{i}\tau^{2})=-U_{ij}\\
a_0^{l}(\pmb{i})\tau^{l}&\to a_0^{\prime l}(\pmb{i})\tau^{l}=(-\mathrm{i}\tau^{2})(a_0^{l}(\pmb{i})\tau^{l})^{\top}(\mathrm{i}\tau^{2})=-a_0^{l}(\pmb{i})\tau^{l}
\end{aligned}
\tag{9.2.44}
\end{equation*}

% ------------------------------------------------------------
% 新增术语（ch9_c，书页 380–392，供术语表收录）：
%   loop variable 圈变量
%   SU(2)-flux operator SU(2) 通量算符
%   base point 基点
%   trivial / collinear / non-collinear SU(2) flux 平凡 / 共线 / 非共线 SU(2) 通量
%   mass term 质量项（ch3_e 已收）
%   uniform RVB state 均匀 RVB 态
%   staggered-flux state 交错通量态
%   d-wave pairing state d 波配对态
%   symmetric spin-liquid state 对称自旋液体态
%   Z₂-symmetric spin-liquid state Z₂ 对称自旋液体态
%   Z₂-gapped spin liquid Z₂ 有能隙自旋液体
%   Z₂-gapless spin liquid Z₂ 无能隙自旋液体
%   Z₂-linear spin liquid Z₂ 线性自旋液体（SU(2)/U(1)-linear state 同式：线性态）
%   sRVB (short-ranged resonating valence bond) state 短程共振价键（sRVB）态
%   Fermi pocket 费米口袋
%   topological soliton 拓扑孤子
%   unfrustrated / frustrated ansatz 非阻挫 / 阻挫拟设
%   invariant group 不变群
%   level (of Chern–Simons theory) level（Chern–Simons 理论的 level，不译）
%   spin–rotation transformation 自旋旋转变换
%   spin-reversal transformation 自旋反转变换
% ------------------------------------------------------------

\end{document}

\begin{document}
\mainmatter
\setcounter{chapter}{9}
\setcounter{section}{3}
% ch10_b 第10章 §10.2尾–§10.3 (书页 453–465 / p468–p480)
% 块界说明：
%   起点——书页 453（p468）顶部 "will be called a string-net operator." 是书页 452
%   （p467，ch10_a 块）末行显示公式 W(C_net)=W(C_1)W(C_2)… 之后半句的延续；
%   ch10_a 已把前半句译文（"……将被称为"）留在其正文最后。本块正文以 %JOIN%
%   起头续译"弦网算符。……"，不设 %--A-TAIL-- 区，亦不重复前半句。
%   终点——书页 465（p480）末句 "The PSGs obtained in Section 10.2 ... agree
%   exactly with the" 延续到书页 466（p481，ch10_c 块）顶部 "fermion PSGs that
%   we obtained above. ..."。按约定，半句译文（"……与"）留在本文件正文最末，
%   未写入任何 TAIL 区；请 ch10_c 以 %JOIN% 起头续译"我们在上面得到的费米子
%   PSG 完全一致。……"，勿重复本块已有内容。
%JOIN%
弦网算符（string-net operator）。只要 $C_i$ 中至少有一个是开弦，态 $|C_1C_2\dots\rangle$ 就是开弦网态。相应的算符 $W(C_{\mathrm{net}})$ 将被称为开弦网算符。若所有 $C_i$ 都是闭合回路，则 $|C_1C_2\dots\rangle$ 就是闭合弦网态，而 $W(C_{\mathrm{net}})$ 是闭合弦网算符。

哈密顿量 $H_J$ 没有弦网凝聚，因为它的基态 $|0\rangle$ 不包含任何弦网，因而 $\langle0|W(C_{\mathrm{net}})|0\rangle=0$。要得到具有闭合弦网凝聚的哈密顿量，我们首先需要找到一个哈密顿量，其基态包含大量任意尺寸的闭合弦网，并且不包含任何开弦（或开弦网）。

让我们先写出一个使闭合弦和闭合弦网不耗费能量、而任何开弦都耗费大量能量的哈密顿量。这样一个哈密顿量可以取如下形式
\begin{equation*}
H_U=-U\sum_{\mathrm{even}}\hat F_{\pmb{i}}
\end{equation*}
其中
\begin{equation*}
\hat F_{\pmb{i}}=\sigma_{\pmb{i}}^{x}\sigma_{\pmb{i}+\pmb{x}}^{y}\sigma_{\pmb{i}+\pmb{x}+\pmb{y}}^{x}\sigma_{\pmb{i}+\pmb{y}}^{y}\tag{10.3.2}
\end{equation*}

我们发现，无弦态 $|0\rangle$ 是 $H_U$ 的基态之一（假定 $U>0$），能量为 $-UN_{\mathrm{site}}$。所有闭合弦态，例如 $W(C_{\mathrm{close}})|0\rangle$，也都是 $H_U$ 的基态，因为 $[H_U,W(C_{\mathrm{close}})]=0$。类似地，$[H_U,W(C_{\mathrm{net}})]=0$ 意味着任何闭合弦态也都是基态。

利用开弦算符 $W(C_{\mathrm{open}})$ 与 $\hat F_{\pmb{i}}$ 之间的对易关系
\begin{align*}
\hat F_{\pmb{i}}W(C_{\mathrm{open}})&=-W(C_{\mathrm{open}})\hat F_{\pmb{i}}, && \text{若 } C_{\mathrm{open}} \text{ 终止于方格 } \pmb{i},\\
\hat F_{\pmb{i}}W(C_{\mathrm{open}})&=+W(C_{\mathrm{open}})\hat F_{\pmb{i}}, && \text{其他情形},
\end{align*}
我们发现，开弦算符在其两个端点处翻转 $\hat F_{\pmb{i}}$ 的符号。开弦态 $W(C_{\mathrm{open}})|0\rangle$ 也是 $\hat F_{\pmb{i}}$ 的本征态，因而是 $H_U$ 的本征态，能量为 $-UN_{\mathrm{site}}+4U$。我们看到，开弦的每个端点耗费能量 $2U$。我们还注意到，闭合弦网的能量不依赖于网中弦的长度。因此，$H_U$ 中的弦没有张力。可以通过在哈密顿量中加入 $H_J$ 来引入弦张力，张力为每格点（或每段）$2J$。我们注意到，任何弦网态 $|C_1C_2\dots\rangle$ 都是 $H_U+H_J$ 的本征态。因此，由 $H_U+H_J$ 描述的弦网不涨落，从而不能凝聚。为使弦发生涨落，我们需要 $g$ 项
\begin{equation*}
H_g=-g\sum_{\pmb{p}}W(C_{\pmb{p}})=-g\sum_{\mathrm{odd}}\hat F_{\pmb{i}}
\end{equation*}
其中 $\pmb{p}$ 标记奇数方格，$C_{\pmb{p}}$ 是围绕方格 $\pmb{p}$ 的闭合弦。当 $H_g$ 作用在弦网态上时，它往弦网中添加一圈弦 $C_{\pmb{p}}$。因此 $g$ 项使弦发生涨落。我们的自旋 1/2 模型的总哈密顿量为
\begin{equation*}
H=H_U+H_J+H_g
\end{equation*}

当 $U\gg g\gg J$ 时，闭合弦几乎没有张力，$H_g$ 诱发的涨落不受张力的限制。因此，基态包含闭合弦的强烈涨落。换言之，基态被任意尺寸的闭合弦所充满。若 $U\gg g$，则闭合弦断开成开弦要耗费过多能量。因此，当 $U\gg g\gg J$ 时，我们自旋 1/2 模型的基态包含闭合弦凝聚（或者更确切地说，闭合弦网凝聚）。下面我们将研究这样一种弦网凝聚态的物理性质。

\subsection*{问题 10.3.1}

(a) 证明各 $\hat F_{\pmb{i}}$ 在 $x$ 与 $y$ 两个方向均为周期性边界条件的格子上满足
\begin{equation*}
\prod_{\pmb{i}}\hat F_{\pmb{i}}=1\tag{10.3.3}
\end{equation*}
式 (10.3.2) 中的负号使上述简单结果成为可能。

(b) 在两方向均为周期性边界条件的偶 $\times$ 偶格子上，证明各 $\hat F_{\pmb{i}}$ 满足更强的条件
\begin{equation*}
\prod_{\pmb{i}=\mathrm{even}}\hat F_{\pmb{i}}=1\tag{10.3.4}
\end{equation*}

\subsection{弦网凝聚与低能有效理论}

\begin{keypoints}\item 弦网凝聚给出规范激发。\end{keypoints}

让我们先讨论 $U\gg g>0$、$J=0$ 的弦网凝聚相。当 $J=0$ 时，模型 $H_U+H_g$ 成为上一节讨论过的模型 (10.2.6)，其中在偶数格点 $V_{\pmb{i}}=U$，在奇数格点 $V_{\pmb{i}}=g$。该模型是精确可解的，因为 $[\hat F_{\pmb{i}},\hat F_{\pmb{j}}]=0$（Kitaev, 2003）。$H_U+H_g$ 的本征态可以由 $\hat F_{\pmb{i}}$ 的共同本征态得到，即 $|\{F_{\pmb{i}}\}\rangle$，其中 $F_{\pmb{i}}$ 是 $\hat F_{\pmb{i}}$ 的本征值。由于 $\hat F_{\pmb{i}}^{2}=1$，$\hat F_{\pmb{i}}$ 的本征值就是 $F_{\pmb{i}}=\pm1$。$|\{F_{\pmb{i}}\}\rangle$ 的能量为 $E=-U\sum_{\pmb{i}=\mathrm{even}}F_{\pmb{i}}-g\sum_{\pmb{i}=\mathrm{odd}}F_{\pmb{i}}$。

注意到，对于尺寸为 $L_x\times L_y$ 的偶 $\times$ 偶周期格子上的有限系统，各 $\hat F_{\pmb{i}}$ 满足约束 $\prod_{\pmb{i}=\mathrm{even}}\hat F_{\pmb{i}}=1$ 与 $\prod_{\pmb{i}=\mathrm{odd}}\hat F_{\pmb{i}}=1$。因此，各 $F_{\pmb{i}}$ 并不独立，只能标记 $2^{L_xL_y}/4$ 种不同的组态。然而，我们自旋 1/2 模型的希尔伯特空间包含 $2^{L_xL_y}$ 个态。为了复现这 $2^{L_xL_y}$ 个态，$\hat F_{\pmb{i}}$ 的共同本征态必定是四重简并的。这与上一节得到的结果一致。

在 $U\gg g$ 的极限下，所有包含开弦的态都将具有量级为 $U$ 的能量。低能态只包含闭合弦网，并且满足
\begin{equation*}
F_{\pmb{i}}\big|_{\pmb{i}=\mathrm{even}}=1
\end{equation*}
不同的低能态由奇数方格上的 $F_{\pmb{i}}$ 标记，即 $F_{\pmb{i}}\big|_{\pmb{i}=\mathrm{odd}}=\pm1$。特别地，基态由下式给出
\begin{equation*}
F_{\pmb{i}}\big|_{\pmb{i}=\mathrm{odd}}=1.
\end{equation*}

这样的基态具有闭合弦凝聚。这是因为所有闭合弦算符 $W(C_{\mathrm{close}})$ 都与 $H_U+H_g$ 对易。于是，$H_U+H_g$ 的基态 $|\Psi_0\rangle$ 是 $W(C_{\mathrm{close}})$ 的本征态，并且满足
\begin{equation*}
\langle\Psi_0|W(C_{\mathrm{close}})|\Psi_0\rangle=1.
\end{equation*}

上式还意味着基态 $|\Psi_0\rangle$ 是各种闭合弦网组态的等权叠加。

基态之上的低能激发可以通过把某些奇数方格上的 $F_{\pmb{i}}$ 从 1 翻转到 $-1$ 得到。若把奇数方格上的 $F_{\pmb{i}}$ 视为 Z$_2$ 规范理论中的通量，我们发现模型的低能部分与 Z$_2$ 格点规范理论完全相同——至少对无限系统如此。这表明我们模型的低能有效理论是 Z$_2$ 格点规范理论。

然而，有人可能会反对这一结果，指出我们模型的低能部分也同样与在每个奇数方格上放一个自旋的伊辛模型相同，因此低能有效理论应当是伊辛模型。我们想指出，尽管对无限系统而言我们模型的低能部分与伊辛模型相同，但对有限系统而言，我们模型的低能部分不同于伊辛模型。例如，在具有周期性边界条件的有限偶 $\times$ 偶格子上，我们模型的基态有四重简并（Kitaev, 2003; Wen, 2003c）。伊辛模型没有这种简并，而 Z$_2$ 规范理论有这种简并。此外，我们的模型包含一种可以认定为 Z$_2$ 电荷（Z$_2$ charge）的激发。因此，我们模型的低能有效理论是 Z$_2$ 格点规范理论，而不是伊辛模型。奇数方格上 $F_{\pmb{i}}=-1$ 的激发可以视为 Z$_2$ 格点规范理论中的 Z$_2$ 涡旋激发。

为理解 Z$_2$ 电荷激发，我们注意到，在弦凝聚态中，弦不耗费能量，并且是不可观测的。把闭合弦算符作用到基态上，得到的仍是基态自身。由于弦不可观测，一段开弦的行为就像两个独立的粒子。开弦的每个端点对应偶数方格上的一个粒子。产生这样一个粒子所需的能量为 $U$ 的量级。现在，让我们考虑这样一个粒子围绕四个最近邻偶数方格的跳跃（见图 10.2）。每一步跳跃由一小段弦算符产生。我们看到，四个跳跃振幅的乘积由这四个偶数方格正中央那个奇数方格上 $\hat F_{\pmb{i}}$ 的本征值给出。这正是电荷与通量之间的关系。因此，若把奇数方格上的 $F_{\pmb{i}}$ 认定为 Z$_2$ 通量，则偶数方格上弦的端点就对应于 Z$_2$ 电荷。由于闭合弦凝聚，开弦的端点不被禁闭，彼此之间只有短程相互作用。因此，Z$_2$ 电荷的行为就像不挂着弦的点粒子。

%FIG{10.2}: {Z$_2$ 电荷围绕四个最近邻偶数方格的一次跳跃。}

\subsection{三种类型的弦与演生费米子}

\begin{keypoints}\item 弦的端点是规范荷。某些弦的端点还可以是费米子。\end{keypoints}

下面我们将研究几种不同类型的弦。为避免混淆，我们把上面讨论的弦称为 T1 弦。T1 弦连接偶数方格（见图 10.3(a)）。

正如 Z$_2$ 电荷一样，一对 Z$_2$ 涡旋也由开弦算符产生。由于 Z$_2$ 涡旋对应于奇数方格上被翻转的 $\hat F_{\pmb{i}}$，产生 Z$_2$ 涡旋的开弦算符同样由式 (10.3.1) 给出，只是现在乘积遍及一条连接\emph{奇数}方格的弦。我们把这样的弦称为 T2 弦。

我们想指出，T2 弦的参考态（即无弦态）不同于 T1 弦的参考态。无 T2 弦态由 $|\tilde{0}\rangle$ 给出，其自旋在偶数格点指向 $y$ 方向、在奇数格点指向 $x$ 方向。由于 T1 弦与 T2 弦有不同的参考态，我们不能同时拥有 T1 弦和 T2 弦的稀薄气体。容易验证，T2 弦算符也与 $H_J+H_g$ 对易。因此，基态 $|\Psi_0\rangle$ 除了 T1 闭合弦的凝聚之外，还有 T2 闭合弦的凝聚。

%FIG{10.3}: {三种类型的弦：(a) T1 弦；(b) T2 弦；(c) T3 弦。}

Z$_2$ 涡旋的跳跃由一段短的 T2 开弦诱导。由于 T2 开弦算符彼此都对易，Z$_2$ 涡旋的行为像玻色子。类似地，Z$_2$ 电荷的行为也像玻色子。然而，T1 开弦算符与 T2 开弦算符不对易。结果是，T1 弦的端点与 T2 弦的端点具有非平凡的互统计。鉴于我们已经证明，让 Z$_2$ 电荷绕 Z$_2$ 涡旋运动一周会产生相位 $\pi$，Z$_2$ 电荷与 Z$_2$ 涡旋具有半子型（semionic）互统计。

T3 弦定义为 T1 弦与 T2 弦的束缚态（见图 10.3(c)）。T3 弦算符由一个 T1 弦算符与一个 T2 弦算符的乘积给出，具有如下形式
\begin{equation*}
W(C)=\prod_n\sigma_{\pmb{i}_n}^{l_n}
\end{equation*}
其中 $C$ 是一条连接相邻链（link）中点的弦，$\pmb{i}_n$ 是弦上的格点。这里，若弦在格点 $\pmb{i}_n$ 处不拐弯，则 $l_n=z$；若弦在格点 $\pmb{i}_n$ 处拐弯，则 $l_n=x$ 或 $y$；若拐弯构成右上角或左下角，则 $l_n=x$；若拐弯构成右下角或左上角，则 $l_n=y$（见图 10.3(c)）。基态也有 T3 闭合弦的凝聚。T3 弦的端点是由一条链两侧两个方格上的一个 Z$_2$ 电荷与一个 Z$_2$ 涡旋形成的束缚态（即链两侧的 $F_{\pmb{i}}=-1$）。这样的束缚态是费米子（见 7.1.2 节与图 7.8）。所以费米子生活在链上。有趣的是，我们模型中的弦网凝聚直接导致 Z$_2$ 规范结构和三种新类型的准粒子，即 Z$_2$ 电荷、Z$_2$ 涡旋和费米子。费米子作为开 T3 弦的端点，从我们这个纯玻色模型中演生出来！

由于 T1 弦的端点是 Z$_2$ 电荷，T1 弦可以视为 Z$_2$“电”通量的弦。类似地，T2 弦可以视为 Z$_2$“磁”通量的弦。

\section{用投影对称群对不同的弦网凝聚分类}

\subsection{四类弦网凝聚}

\begin{keypoints}\item 不同的弦网凝聚给出不同的量子有序相。\end{keypoints}

我们已经看到，当 $U>0$、$g>0$、$J=0$ 时，我们模型 $H_g+H_U+H_J$ 的基态由下式给出
\begin{equation*}
F_{\pmb{i}}\big|_{\pmb{i}=\mathrm{even}}=1,\qquad F_{\pmb{i}}\big|_{\pmb{i}=\mathrm{odd}}=1.
\end{equation*}
我们将把这样的相称为 Z$_2$ 相，以强调其低能 Z$_2$ 规范结构。在 Z$_2$ 相中，T1 弦算符 $W_1(C_1)$ 与 T2 弦算符 $W_2(C_2)$ 有如下期望值：
\begin{equation*}
\langle W_1(C_1)\rangle=1,\qquad \langle W_2(C_2)\rangle=1
\end{equation*}

当 $U>0$、$g<0$、$J=0$ 时，基态由下式给出
\begin{equation*}
F_{\pmb{i}}\big|_{\pmb{i}=\mathrm{even}}=1,\qquad F_{\pmb{i}}\big|_{\pmb{i}=\mathrm{odd}}=-1.
\end{equation*}
我们看到每个奇数方格都有 $\pi$ 通量。我们将把这样的相称为 Z$_2$ 通量相。T1 弦算符与 T2 弦算符有如下期望值：
\begin{equation*}
\langle W_1(C_1)\rangle=(-)^{N_{\mathrm{odd}}},\qquad \langle W_2(C_2)\rangle=1
\end{equation*}
其中 $N_{\mathrm{odd}}$ 是被 T1 弦 $C_1$ 包围的奇数方格的数目。

当 $U<0$、$g>0$、$J=0$ 时，基态由下式给出
\begin{equation*}
F_{\pmb{i}}\big|_{\pmb{i}=\mathrm{even}}=-1,\qquad F_{\pmb{i}}\big|_{\pmb{i}=\mathrm{odd}}=1.
\end{equation*}
每个偶数方格上有一个 Z$_2$ 电荷。我们将把这样的相称为 Z$_2$ 电荷相。T1 弦算符与 T2 弦算符有如下期望值：
\begin{equation*}
\langle W_1(C_1)\rangle=1,\qquad \langle W_2(C_2)\rangle=(-)^{N_{\mathrm{even}}}
\end{equation*}
其中 $N_{\mathrm{even}}$ 是被 T2 弦 $C_2$ 包围的偶数方格的数目。注意，Z$_2$ 通量相与 Z$_2$ 电荷相仅相差一个格子平移，本质上是同一个相。

当 $U<0$、$g<0$、$J=0$ 时，基态变为
\begin{equation*}
F_{\pmb{i}}\big|_{\pmb{i}=\mathrm{even}}=-1,\qquad F_{\pmb{i}}\big|_{\pmb{i}=\mathrm{odd}}=-1.
\end{equation*}
每个偶数方格上有一个 Z$_2$ 电荷，且每个奇数方格有 $\pi$ 通量。我们将把这样的相称为 Z$_2$ 通量--电荷相。T1 弦算符与 T2 弦算符有如下期望值：
\begin{equation*}
\langle W_1(C_1)\rangle=(-)^{N_{\mathrm{odd}}},\qquad \langle W_2(C_2)\rangle=(-)^{N_{\mathrm{even}}}
\end{equation*}

当 $U=g=0$、$J\neq0$ 时，模型同样精确可解。基态是一个简单的自旋极化态（见图 10.1）。

从 $H_U+H_g+H_J$ 模型的这些精确可解极限出发，我们提出一幅相图，草图示于图 10.4。该相图包含四个不同的弦网凝聚相和一个无弦凝聚的相。所有的相都有相同的对称性，仅凭它们不同的量子序相互区分。

\subsection{投影对称群与凝聚弦的端点}

\begin{keypoints}\item PSG 所描述的投影对称性，正是凝聚弦端点有效理论的对称性。\end{keypoints}

从各不相同的 $\langle W_1(C_1)\rangle$ 与 $\langle W_2(C_2)\rangle$ 我们看到，前四个相具有不同的弦网凝聚。然而，它们都有相同的对称性。这就引出一个问题：在没有对称性破缺的情况下，我们怎么知道上述四个相的确是不同的相？我们怎么知道不可能在不经历相变的情况下把一个弦网凝聚态变成另一个弦网凝聚态？

%FIG{10.4}: {$H=H_U+H_g+H_J$ 模型的建议相图。这里假定 $J$ 为正。四个弦网凝聚相由一对 PSG（$PSG_{\mathrm{charge}}$，$PSG_{\mathrm{vortex}}$）表征。SP 标记自旋极化相。}

下面我们将证明，不同的弦网凝聚可以用不同的 PSG 来描述（正如不同的对称性破缺序可以用基态的不同对称群来描述）。在第 9 章中，不同的量子序是通过它们不同的 PSG 引入的（Wen, 2002b,c）。那里证明了 PSG 是量子相的普适性质，只有相变才能改变 PSG。因此，不同的弦网凝聚态由不同的 PSG 表征，这一事实表明这些不同的弦网凝聚态属于不同的量子相。

弦网凝聚与 PSG 之间的联系也使我们能够把弦网凝聚与第 9 章引入的量子序联系起来。事实上，第 9 章讨论的投影构造可以视为构造弦网凝聚态的一种方法。

为看清 PSG 与弦网凝聚之间的联系，我们注意到，当闭合弦网凝聚时，开弦的端点表现得像独立的粒子。让我们考虑描述 T1 弦两个端点的双粒子态 $|\pmb{p}_1\pmb{p}_2\rangle$。注意，T1 弦的端点（因而 Z$_2$ 电荷）只生活在偶数方格上。因此，$\pmb{p}_1$ 与 $\pmb{p}_2$ 只标记偶数方格。对于我们的模型 $H_U+H_g$，$|\pmb{p}_1\pmb{p}_2\rangle$ 是能量本征态，而且 Z$_2$ 电荷不跳跃。这里我们想往哈密顿量中加入一项
\begin{equation*}
H_t=t\sum_{\pmb{i}}(\sigma_{\pmb{i}}^{x}+\sigma_{\pmb{i}}^{y})+t'\sum_{\pmb{i}}\sigma_{\pmb{i}}^{z}
\end{equation*}
$t$ 项 $t\sum_{\pmb{i}}(\sigma_{\pmb{i}}^{x}+\sigma_{\pmb{i}}^{y})$ 使 Z$_2$ 电荷在偶数方格之间跳跃，跳跃幅为 $t$ 的量级。两个 Z$_2$ 电荷的动力学由双粒子希尔伯特空间中的如下有效哈密顿量描述：
\begin{equation*}
H=H(\pmb{p}_1)+H(\pmb{p}_2)
\end{equation*}
其中 $H(\pmb{p}_1)$ 描述第一个粒子 $\pmb{p}_1$ 的跳跃，$H(\pmb{p}_2)$ 描述第二个粒子 $\pmb{p}_2$ 的跳跃。现在我们可以定义弦网凝聚态的 PSG 了：PSG 就是跳跃哈密顿量 $H(\pmb{p})$ 的对称群。

由于底层模型 $H_U+H_g+H_t$ 具有平移对称性，我们可能天真地预期 Z$_2$ 电荷的跳跃哈密顿量 $H(\pmb{p})$ 在 $\pmb{x}+\pmb{y}$ 与 $\pmb{x}-\pmb{y}$ 方向也具有平移对称性：
\begin{align*}
H(\pmb{p})&=T_{xy}^{\dagger}H(\pmb{p})T_{xy}, & T_{xy}|\pmb{p}\rangle&=|\pmb{p}+\pmb{x}+\pmb{y}\rangle\\
H(\pmb{p})&=T_{x\bar{y}}^{\dagger}H(\pmb{p})T_{x\bar{y}}, & T_{x\bar{y}}|\pmb{p}\rangle&=|\pmb{p}+\pmb{x}-\pmb{y}\rangle\tag{10.4.1}
\end{align*}
如果上式成立，则意味着 PSG 与平移对称群相同。

结果表明，式 (10.4.1) 过强。即使 $H(\pmb{p})$ 不满足式 (10.4.1)，底层自旋模型仍然可以具有平移对称性。然而，$H(\pmb{p})$ 可能的对称群（即各个 PSG）受到底层自旋模型平移对称性的很强约束。下面我们将解释，为什么 PSG 可以不同于物理自旋模型的对称群，以及为与底层自旋模型的平移对称性相一致，PSG 必须满足什么条件。

我们注意到，弦总有两个端点。因此，一个物理态总是拥有偶数个 Z$_2$ 电荷。平移作用在双粒子态上为
\begin{align*}
T_{xy}^{(2)}|\pmb{p}_1,\pmb{p}_2\rangle&=|\pmb{p}_1+\pmb{x}+\pmb{y},\ \pmb{p}_2+\pmb{x}+\pmb{y}\rangle\\
T_{x\bar{y}}^{(2)}|\pmb{p}_1,\pmb{p}_2\rangle&=|\pmb{p}_1+\pmb{x}-\pmb{y},\ \pmb{p}_2+\pmb{x}-\pmb{y}\rangle
\end{align*}
这里 $T_{xy}^{(2)}$ 与 $T_{x\bar{y}}^{(2)}$ 满足平移代数
\begin{equation*}
T_{xy}^{(2)}T_{x\bar{y}}^{(2)}=T_{x\bar{y}}^{(2)}T_{xy}^{(2)}\tag{10.4.2}
\end{equation*}

我们注意到，$T_{xy}^{(2)}$ 与 $T_{x\bar{y}}^{(2)}$ 是单粒子态上平移算符的直积。因此，在某种意义上，单粒子平移是双粒子平移的平方根。双粒子平移代数对单粒子平移代数施加了很强的约束。

单粒子平移最一般的形式由 $T_{xy}G_{xy}$ 与 $T_{x\bar{y}}G_{x\bar{y}}$ 给出，其中算符 $T_{xy,x\bar{y}}$ 与 $G_{xy,x\bar{y}}$ 的作用定义为
\begin{align*}
T_{xy}|\pmb{p}\rangle&=|\pmb{p}+\pmb{x}+\pmb{y}\rangle, & T_{x\bar{y}}|\pmb{p}\rangle&=|\pmb{p}+\pmb{x}-\pmb{y}\rangle,\\
G_{xy}|\pmb{p}\rangle&=\mathrm{e}^{\phi_{xy}(\pmb{p})}|\pmb{p}\rangle, & G_{x\bar{y}}|\pmb{p}\rangle&=\mathrm{e}^{\phi_{x\bar{y}}(\pmb{p})}|\pmb{p}\rangle.
\end{align*}
为了使直积 $T_{xy}^{(2)}=T_{xy}G_{xy}\otimes T_{xy}G_{xy}$ 与 $T_{x\bar{y}}^{(2)}=T_{x\bar{y}}G_{x\bar{y}}\otimes T_{x\bar{y}}G_{x\bar{y}}$ 复现平移代数 (10.4.2)，我们只要求 $T_{xy}G_{xy}$ 与 $T_{x\bar{y}}G_{x\bar{y}}$ 满足
\begin{equation*}
T_{xy}G_{xy}T_{x\bar{y}}G_{x\bar{y}}=T_{x\bar{y}}G_{x\bar{y}}T_{xy}G_{xy}\tag{10.4.3}
\end{equation*}
或
\begin{equation*}
T_{xy}G_{xy}T_{x\bar{y}}G_{x\bar{y}}=-T_{x\bar{y}}G_{x\bar{y}}T_{xy}G_{xy}\tag{10.4.4}
\end{equation*}

算符 $T_{xy}G_{xy}$ 与 $T_{x\bar{y}}G_{x\bar{y}}$ 生成一个群，这个群正是 PSG。两个不同的代数 (10.4.3) 与 (10.4.4) 生成两个不同的 PSG，二者都与作用在双粒子态上的平移群相一致。我们把由式 (10.4.3) 生成的 PSG 称为 Z2a PSG，把由式 (10.4.4) 生成的 PSG 称为 Z2b PSG。

让我们回顾一下 PSG 的一般定义。PSG 是一个群，它是对称群（SG）的扩张，即 PSG 包含一个正规子群（称为不变规范群或 IGG），使得
\begin{equation*}
PSG/IGG=SG
\end{equation*}

在我们的情形中，SG 是平移群 $SG=\{1,T_{xy}^{(2)},T_{x\bar{y}}^{(2)},\dots\}$。IGG 由单粒子态上满足 $G_0\otimes G_0=1$ 的变换 $G_0$ 构成。我们发现 IGG 由下式生成：
\begin{equation*}
G_0|\pmb{p}\rangle=-|\pmb{p}\rangle
\end{equation*}

单粒子态上的这三个变换 $G_0$、$T_{xy}G_{xy}$ 与 $T_{x\bar{y}}G_{x\bar{y}}$ 生成 Z2a 或 Z2b PSG。

我们现在看到，底层平移对称性并不要求单粒子跳跃哈密顿量 $H(\pmb{p})$ 具有平移对称性，只要求 $H(\pmb{p})$ 在 Z2a PSG 或 Z2b PSG 下不变。当 $H(\pmb{p})$ 在 Z2a PSG 下不变时，跳跃哈密顿量具有通常的平移对称性。当 $H(\pmb{p})$ 在 Z2b PSG 下不变时，跳跃哈密顿量具有磁平移对称性，描述磁场中每个奇数方格带 $\pi$ 通量的跳跃。

\subsection{投影对称群对不同的弦网凝聚的分类}

\begin{keypoints}\item 不同的弦网凝聚有不同的 PSG。由于 PSG 是量子相的普适性质，我们发现不同的弦网凝聚对应于不同的相。\end{keypoints}

在了解了弦端点的跳跃哈密顿量可能具有的 PSG 之后，我们现在可以计算实际的 PSG 了。让我们考虑模型 $H_U+H_g+H_t$ 的两个基态：一个具有 $F_{\pmb{i}}\big|_{\pmb{i}=\mathrm{odd}}=1$（对应 $g>0$），另一个具有 $F_{\pmb{i}}\big|_{\pmb{i}=\mathrm{odd}}=-1$（对应 $g<0$）。两个基态在 $\pmb{x}+\pmb{y}$ 与 $\pmb{x}-\pmb{y}$ 方向都有相同的平移对称性。然而，相应的单粒子跳跃哈密顿量 $H(\pmb{p})$ 却有不同的对称性。对于 $F_{\pmb{i}}\big|_{\pmb{i}=\mathrm{odd}}=1$ 的态，奇数方格没有通量，$H(\pmb{p})$ 具有通常的平移对称性，在 Z2a PSG 下不变。对于 $F_{\pmb{i}}\big|_{\pmb{i}=\mathrm{odd}}=-1$ 的态（译注：原文此处印作“$F_{\pmb{i}}|_{\pmb{i}=\mathrm{odd}}=1$”，按上下文应为 $-1$），奇数方格有 $\pi$ 通量，$H(\pmb{p})$ 具有磁平移对称性，其 PSG 为 Z2b PSG。因此，$F_{\pmb{i}}\big|_{\pmb{i}=\mathrm{odd}}=1$ 的态与 $F_{\pmb{i}}\big|_{\pmb{i}=\mathrm{odd}}=-1$ 的态尽管对称性相同，却有不同的序。两个态中不同的量子序可以用它们不同的 PSG 来表征。

如 10.4.1 小节所述，上述两个态有不同的弦网凝聚。对 $\hat F_{\pmb{i}}\big|_{\pmb{i}=\mathrm{odd}}=1$ 的态，T1 闭合弦算符的平均值为 $\langle W(C)\rangle=1$；而对 $\hat F_{\pmb{i}}\big|_{\pmb{i}=\mathrm{odd}}=-1$ 的态，T1 闭合弦算符的平均值为 $\langle W(C)\rangle=(-)^{N_{\mathrm{odd}}(C)}$，其中 $N_{\mathrm{odd}}(C)$ 是被 T1 闭合弦 $C$ 包围的奇数方格的数目。因此，我们也可以说，不同的弦网凝聚可以用不同的 PSG 来表征。

在上面的讨论中，我们只在一些特定的精确可解模型中表明了 PSG 可以表征不同的弦网凝聚态。事实上，在一般模型中，不同的 PSG 确实表征不同的量子相。这是因为，正如 9.9.1 小节所证明的，PSG 是相的普适性质（Wen, 2002c）。弦网凝聚态的 PSG 对物理模型的任何不改变模型对称性的局域微扰都是稳健的。

上述讨论同样适用于 Z$_2$ 涡旋与 T2 弦。因此，我们模型中的量子序由一对 PSG（$PSG_{\mathrm{charge}}$，$PSG_{\mathrm{vortex}}$）描述，一个对应 Z$_2$ 电荷，一个对应 Z$_2$ 涡旋。PSG 对（$PSG_{\mathrm{charge}}$，$PSG_{\mathrm{vortex}}$）使我们能够区分模型 $H=H_U+H_g+H_t$ 的四个不同的弦网凝聚态（见图 10.5）。

\subsection{T3 弦端点的投影对称群}

\begin{keypoints}\item 对于具有几种不同类型弦凝聚的态，不同凝聚弦的端点可以有不同的投影对称性和不同的 PSG。\end{keypoints}

%FIG{10.5}: {$H=H_U+H_g+H_t$ 模型的建议相图。这里假定 $t=t'$ 为正。四个弦网凝聚相由 PSG 对（$PSG_{\mathrm{charge}}$，$PSG_{\mathrm{vortex}}$）表征。SP 标记均匀自旋极化相。所有的相都有相同的对称性，并且没有任何相变改变对称性。}

\begin{keypoints}\item 9.4.2 小节引入的 PSG 对应于 T3 弦的 PSG。\end{keypoints}

在本小节中，我们将假设模型中 $U=g=V$：
\begin{equation*}
H_U+H_g+H_t=H_t-V\sum_{\pmb{i}}\hat F_{\pmb{i}}\tag{10.4.5}
\end{equation*}

新模型具有由 $\pmb{i}\to\pmb{i}+\pmb{x}$ 与 $\pmb{i}\to\pmb{i}+\pmb{y}$ 生成的更大的平移对称性（见图 10.5）。注意到，当 $H_t=0$ 时（译注：原文此处印作“$H_l=0$”，应为 $H_t$），新模型与 10.2.1 小节研究过的自旋哈密顿量 (10.2.6) 相同。凝聚的 T1 弦与 T2 弦的 PSG 已在上一小节研究过，这里我们想讨论 T3 弦的 PSG。由于 T3 弦的端点生活在链上，相应的单粒子跳跃哈密顿量 $H_f(\pmb{l})$ 描述费米子在链之间的跳跃。显然，$H_f(\pmb{l})$ 的对称群（即 PSG）可以不同于 $H(\pmb{p})$ 的对称群。

让我们把费米子跳跃 $\pmb{l}\to\pmb{l}+\pmb{x}$ 定义为两步跳跃 $\pmb{l}\to\pmb{l}+\frac{\pmb{x}}{2}-\frac{\pmb{y}}{2}\to\pmb{l}+\pmb{x}$ 的组合（见图 10.6）。跳跃 $\pmb{l}\to\pmb{l}+\frac{\pmb{x}}{2}-\frac{\pmb{y}}{2}$ 由 $\sigma_{\pmb{l}+\frac{\pmb{x}}{2}}^{x}$ 产生，跳跃 $\pmb{l}+\frac{\pmb{x}}{2}-\frac{\pmb{y}}{2}\to\pmb{l}+\pmb{x}$ 由 $\sigma_{\pmb{l}+\frac{\pmb{x}}{2}}^{y}$ 产生。因此，跳跃 $\pmb{l}\to\pmb{l}+\pmb{x}$ 由 $\sigma_{\pmb{l}+\frac{\pmb{x}}{2}}^{y}\sigma_{\pmb{l}+\frac{\pmb{x}}{2}}^{x}$ 产生。类似地，我们把费米子跳跃 $\pmb{l}\to\pmb{l}+\pmb{y}$ 定义为两步跳跃 $\pmb{l}\to\pmb{l}+\frac{\pmb{x}}{2}+\frac{\pmb{y}}{2}\to\pmb{l}+\pmb{y}$，它由 $\sigma_{\pmb{l}+\frac{\pmb{x}}{2}+\pmb{y}}^{x}\sigma_{\pmb{l}+\frac{\pmb{x}}{2}}^{y}$ 产生。
%FIG{10.6}: {费米子绕一个方格和绕一个格点的跳跃。}
在这样的定义下，费米子绕方格 $\pmb{l}\to\pmb{l}+\pmb{x}\to\pmb{l}+\pmb{x}+\pmb{y}\to\pmb{l}+\pmb{y}\to\pmb{l}$ 跳跃一周由下式产生（见图 10.6）：$(\sigma_1^{y}\sigma_4^{x})(\sigma_4^{x}\sigma_4^{y})(\sigma_3^{x}\sigma_2^{y})(\sigma_1^{y}\sigma_1^{x})=\sigma_4^{y}\sigma_3^{x}\sigma_2^{y}\sigma_1^{x}=\hat F_1$，或者 $(\sigma_5^{z})(\sigma_5^{y}\sigma_8^{x})(\sigma^{z})(\sigma_7^{x}\sigma_6^{y})=-\sigma_5^{x}\sigma_8^{y}\sigma_7^{x}\sigma_6^{y}=-\hat F_5$。基态具有 $F_{\pmb{i}}=\mathrm{sgn}(V)$。然而，由于 T3 弦终止在链 $5\to6$ 上，我们有 $-\hat F_5=\hat F_{\pmb{i}}=\mathrm{sgn}(V)$；另一方面，我们有 $\hat F_1=\hat F_{\pmb{i}}=\mathrm{sgn}(V)$。因此，费米子绕 1234 与 5678 两个方格跳跃的总振幅由式 (10.4.5) 中 $V$ 的符号给出。当 $\mathrm{sgn}(V)=1$ 时，费米子看不到通量；而当 $\mathrm{sgn}(V)=-1$ 时，费米子看到每个方格有 $\pi$ 通量。结果是，$H_f(\pmb{l})$ 平移对称性的生成元 $(T_xG_x,\,T_yG_y)$ 满足磁平移代数（见式 (7.2.3)）
\begin{equation*}
(T_yG_y)^{-1}(T_xG_x)^{-1}T_yG_yT_xG_x=\mathrm{sgn}(V).\tag{10.4.6}
\end{equation*}

除了在 $(T_xG_x,\,T_yG_y)$ 下不变之外，$H_f(\pmb{l})$ 在 $G_0$ 下也不变：
\begin{equation*}
G_0|\pmb{l}\rangle=-|\pmb{l}\rangle
\end{equation*}
于是，$(G_0,\,T_xG_x,\,T_yG_y)$ 生成 $H_f(\pmb{l})$ 的对称群——费米子 PSG。

当 $\mathrm{sgn}(V)=1$ 时，$T_xG_x$ 与 $T_yG_y$ 满足平移代数，相应的费米子 PSG 是 Z2A PSG。当 $\mathrm{sgn}(V)=-1$ 时，$T_xG_x$ 与 $T_yG_y$ 满足每个方格带 $\pi$ 通量的磁平移代数，相应的费米子 PSG 是 Z2B PSG。我们看到，基态的量子序也可以用费米子 PSG 来表征。$V<0$ 基态的量子序由 Z2A PSG 表征，而 $V>0$ 基态的量子序由 Z2B PSG 表征。

在 10.2 节中，自旋 1/2 模型 (10.4.5)（取 $t=t'=0$）曾通过把自旋拆成两个费米子、用从属玻色子方法求解。那里证明了，$V>0$ 与 $V<0$ 两种态的费米子跳跃哈密顿量具有不同的对称性，或者说在不同的 PSG 下不变。不同的 PSG 意味着 $V>0$ 与 $V<0$ 基态中不同的量子序。10.2 节中针对 $V>0$ 与 $V<0$ 相得到的 PSG 与

与我们在上面得到的费米子 PSG 完全一致。这个例子表明，9.4 节中为描述量子序而引入的各个 PSG，正是凝聚弦端点跳跃哈密顿量的对称群。
\end{document}
